\documentclass[10pt,a4paper]{amsart}
\usepackage{amsmath,amssymb,amsthm,graphicx}
\usepackage{verbatim}
\newenvironment{reference}{\comment}{\endcomment}
%\newenvironment{reference}{}{}
\newenvironment{slogan}{\comment}{\endcomment}
\newenvironment{history}{\comment}{\endcomment}

% For commutative diagrams we use Xy-pic
\usepackage[all]{xy}

% We use 2cell for 2-commutative diagrams.
\xyoption{2cell}
% Theorem environments.
%
\theoremstyle{plain}
\newtheorem{theorem}[subsection]{Theorem}
\newtheorem{proposition}[subsection]{Proposition}
\newtheorem{lemma}[subsection]{Lemma}

\theoremstyle{definition}
\newtheorem{definition}[subsection]{Definition}
\newtheorem{example}[subsection]{Example}
\newtheorem{exercise}[subsection]{Exercise}
\newtheorem{situation}[subsection]{Situation}

\theoremstyle{remark}
\newtheorem{remark}[subsection]{Remark}
\newtheorem{remarks}[subsection]{Remarks}

\numberwithin{equation}{subsection}

% Macros
%
\def\lim{\mathop{\mathrm{lim}}\nolimits}
\def\colim{\mathop{\mathrm{colim}}\nolimits}
\def\Spec{\mathop{\mathrm{Spec}}}
\def\Hom{\mathop{\mathrm{Hom}}\nolimits}
\def\Ext{\mathop{\mathrm{Ext}}\nolimits}
\def\SheafHom{\mathop{\mathcal{H}\!\mathit{om}}\nolimits}
\def\SheafExt{\mathop{\mathcal{E}\!\mathit{xt}}\nolimits}
\def\Sch{\mathit{Sch}}
\def\Mor{\mathop{\mathrm{Mor}}\nolimits}
\def\Ob{\mathop{\mathrm{Ob}}\nolimits}
\def\Sh{\mathop{\mathit{Sh}}\nolimits}
\def\NL{\mathop{N\!L}\nolimits}
\def\CH{\mathop{\mathrm{CH}}\nolimits}
\def\proetale{{pro\text{-}\acute{e}tale}}
\def\etale{{\acute{e}tale}}
\def\QCoh{\mathit{QCoh}}
\def\Ker{\mathop{\mathrm{Ker}}}
\def\Im{\mathop{\mathrm{Im}}}
\def\Coker{\mathop{\mathrm{Coker}}}
\def\Coim{\mathop{\mathrm{Coim}}}

% Boxtimes
%
\DeclareMathSymbol{\boxtimes}{\mathbin}{AMSa}{"02}

%
% Macros for moduli stacks/spaces
%
\def\QCohstack{\mathcal{QC}\!\mathit{oh}}
\def\Cohstack{\mathcal{C}\!\mathit{oh}}
\def\Spacesstack{\mathcal{S}\!\mathit{paces}}
\def\Quotfunctor{\mathrm{Quot}}
\def\Hilbfunctor{\mathrm{Hilb}}
\def\Curvesstack{\mathcal{C}\!\mathit{urves}}
\def\Polarizedstack{\mathcal{P}\!\mathit{olarized}}
\def\Complexesstack{\mathcal{C}\!\mathit{omplexes}}
% \Pic is the operator that assigns to X its picard group, usage \Pic(X)
% \Picardstack_{X/B} denotes the Picard stack of X over B
% \Picardfunctor_{X/B} denotes the Picard functor of X over B
\def\Pic{\mathop{\mathrm{Pic}}\nolimits}
\def\Picardstack{\mathcal{P}\!\mathit{ic}}
\def\Picardfunctor{\mathrm{Pic}}
\def\Deformationcategory{\mathcal{D}\!\mathit{ef}}
\usepackage{cleveref}
\crefname{originalsection}{original section}{original sections}
\Crefname{originalsection}{Original section}{Original sections}
\makeatletter
\newlabel{homology-section-functors}{{010M}{source}{}{sourcepointer-1}{}}
\newlabel{homology-section-functors@cref}{{[originalsection][0][]010M}{[1][0][]source}{}{}{}}
\makeatother
\title{Brown representability in a Grothendieck abelian category}
\author{The Stacks Project Authors}
\date{stacks-project:injectives.tex:local-snapshot-20261012; GNU FDL 1.2 or later}
\begin{document}
\maketitle
\noindent\textit{Editorial note.} All delivered mathematical text is the author's own, copied from the pinned Stacks Project chapter file sites-cohomology.tex; no statement or proof step was rewritten, reordered or removed.\par\medskip
\noindent\textit{Editorial note.} The proof is the author's complete argument; the single omitted step is the sentence 'This construction gives an inverse to the map ... (details omitted).', a routine verification that the map just constructed is inverse to a displayed map. It is preserved verbatim, and the whole construction leading to it is written out in the counted span.\par\medskip
\noindent\textit{Editorial note.} Every labelled reference occurring inside the counted spans is carried in assembly.external\_references as a non-load-bearing navigation pointer (same-chapter ones at their definition site, cross-chapter ones at the invocation inside the counted span), so no load-bearing mathematics is replaced by a stub. No bibliography citation occurs inside the counted spans.\par\medskip
\noindent\textit{Editorial note.} The same-name check for this card scanned every staged Stacks chapter of this repository -- all sources/stacks-*/pinned-original-source.tex files, the card's own chapter excluded -- and the labels found are recorded in dedup.same\_name\_cross\_chapter on the card.\par\medskip
\noindent\textit{Editorial note.} The preamble is the chapter's own retained minimal preamble (sources/stacks-sites-cohomology/retained-preamble.tex), a verbatim excerpt of the native preamble: the Xy-pic setup and the native theorem declarations with the mathematical macros.\par\medskip
\noindent\textit{Allowed standard prerequisites (editorial).}\begin{itemize}
\item The functor represented by an object. Same chapter, the discussion of the functor represented by an object and its universal element. Supplies the description of a representable functor and the universal element used in the construction.
\item The small object argument in the form used here. Same chapter, Lemma proposition-modules-are-small and the surrounding discussion of smallness with respect to injections. Supplies that the objects of a Grothendieck abelian category are small with respect to injections, which produces the subobject generated by the image of a given element.
\item Colimits and their universal property. Same chapter, the universal property of colimits in the category. Supplies the direction in which representable functors commute with colimits and the description of colimits used to build the representing object.
\end{itemize}\medskip
\ifdefined\hypertarget\hypertarget{sourcepointer-1}{}\fi
\noindent\textit{Source reference.} Original section 010M: sources/stacks-homology/pinned-original-source.tex, source lines 1680--1680. Navigation pointer invoked at proof line 2140; the counted span contains the pointer only and the cited result is not reproduced inside it, so the pointer is registered as a non-load-bearing navigation stub in the established form for Stacks material. No bibliography citation is involved.\par\medskip
\begin{lemma}
\label{lemma-grothendieck-brown}
Let $\mathcal{A}$ be a Grothendieck abelian category.
Let $F : \mathcal{A}^{opp} \to \textit{Sets}$ be a functor.
Then $F$ is representable if and only if $F$ commutes with colimits, i.e.,
$$
F(\colim_i N_i) = \lim F(N_i)
$$
for any diagram $\mathcal{I} \to \mathcal{A}$, $i \in \mathcal{I}$.
\end{lemma}

\begin{proof}
If $F$ is representable, then it commutes with colimits by definition
of colimits.

\medskip\noindent
Assume that $F$ commutes with colimits. Then $F(M \oplus N) = F(M) \times F(N)$
and we can use this to define a group structure on $F(M)$. Hence we get
$F : \mathcal{A} \to \textit{Ab}$ which is additive and right exact, i.e.,
transforms a short exact sequence $0 \to K \to L \to M \to 0$ into an exact
sequence $F(K) \leftarrow F(L) \leftarrow F(M) \leftarrow 0$ (compare with
Homology, Section \ref{homology-section-functors}).

\medskip\noindent
Let $U$ be a generator for $\mathcal{A}$. Set $A = \bigoplus_{s \in F(U)} U$.
Let $s_{univ} = (s)_{s \in F(U)} \in F(A) = \prod_{s \in F(U)} F(U)$. Let
$A' \subset A$ be the largest subobject such that $s_{univ}$ restricts to
zero on $A'$. This exists because $\mathcal{A}$ is a Grothendieck category
and because $F$ commutes with colimits. Because $F$ commutes with colimits
there exists a unique element $\overline{s}_{univ} \in F(A/A')$ which
maps to $s_{univ}$ in $F(A)$. We claim that $A/A'$ represents $F$, in
other words, the Yoneda map
$$
\overline{s}_{univ} : h_{A/A'} \longrightarrow F
$$
is an isomorphism. Let $M \in \Ob(\mathcal{A})$ and $s \in F(M)$. Consider
the surjection
$$
c_M :
A_M = \bigoplus\nolimits_{\varphi \in \Hom_\mathcal{A}(U, M)} U
\longrightarrow
M.
$$
This gives $F(c_M)(s) = (s_\varphi) \in \prod_\varphi F(U)$.
Consider the map
$$
\psi :
A_M = \bigoplus\nolimits_{\varphi \in \Hom_\mathcal{A}(U, M)} U
\longrightarrow
\bigoplus\nolimits_{s \in F(U)} U = A
$$
which maps the summand corresponding to $\varphi$ to the summand
corresponding to $s_\varphi$ by the identity map on $U$. Then $s_{univ}$
maps to $(s_\varphi)_\varphi$ by construction.
in other words the right square in the diagram
$$
\xymatrix{
A' \ar[r] &
A \ar@{..>}[r]_{s_{univ}} & F \\
K \ar[r] \ar[u]^{?} & A_M \ar[u]^\psi \ar[r] &
M \ar@{..>}[u]_s
}
$$
commutes. Let $K = \Ker(A_M \to M)$. Since $s$ restricts to zero
on $K$ we see that $\psi(K) \subset A'$ by definition of $A'$. Hence there
is an induced morphism $M \to A/A'$. This construction gives an inverse
to the map $h_{A/A'}(M) \to F(M)$ (details omitted).
\end{proof}

\end{document}
